\section{3D 生成模型}

\subsection{3D 形状生成的早期方法}

早期的 3D 生成模型直接在 3D 数据上训练 GAN 或 VAE：

\begin{figure}[H]
\centering
\includegraphics[width=0.95\textwidth]{slides-images/slide-046.jpg}
\caption{3D GAN：在体素/隐式表示上训练生成对抗网络，生成新的 3D 形状\protect\footnotemark}
\end{figure}
\footnotetext{Slides 第46页。}

\begin{knowledgebox}{3D 生成的局限}
早期方法受限于 3D 数据量少，生成效果有限且局限于特定类别（如椅子、汽车）。突破的关键在于利用 2D 图像/视频基础模型的先验知识来辅助 3D 生成。
\end{knowledgebox}

\subsection{利用 2D 先验的 3D 生成}

通过可微渲染，可以将 3D 生成与 2D 判别器/损失函数连接起来，利用大规模 2D 数据训练的模型来指导 3D 生成。

\begin{figure}[H]
\centering
\includegraphics[width=0.95\textwidth]{slides-images/slide-049.jpg}
\caption{结合 NeRF 与 GAN 的 3D 生成：生成 NeRF 表示的场景，通过可微渲染得到 2D 图像，用 2D 判别器提供梯度\protect\footnotemark}
\end{figure}
\footnotetext{Slides 第49页。}

\subsection{基于扩散模型的 3D 生成}

近年来，Score Distillation Sampling（SDS）等技术使得利用预训练的 2D 扩散模型（如 Stable Diffusion）来指导 3D 生成成为可能。代表性工作包括 DreamFusion、Magic3D 等。

\begin{figure}[H]
\centering
\includegraphics[width=0.95\textwidth]{slides-images/slide-053.jpg}
\caption{基于文本的 3D 生成：通过 SDS 损失将 2D 扩散模型的知识蒸馏到 3D NeRF 中\protect\footnotemark}
\end{figure}
\footnotetext{Slides 第53页。}

\begin{figure}[H]
\centering
\includegraphics[width=0.95\textwidth]{slides-images/slide-056.jpg}
\caption{文本到 3D 的生成结果：输入文本描述，输出可从任意视角渲染的 3D 模型\protect\footnotemark}
\end{figure}
\footnotetext{Slides 第56页。}

\subsection{前沿方向：大规模 3D 生成}

最新的研究方向包括：
\begin{itemize}
    \item \textbf{单图像到 3D}：从一张图片重建完整 3D 模型
    \item \textbf{文本到 3D}：从文字描述直接生成 3D 物体
    \item \textbf{视频到 4D}：从视频重建动态 3D 场景
    \item \textbf{3D 基础模型}：训练能泛化到多类别的统一 3D 生成模型
\end{itemize}

\begin{figure}[H]
\centering
\includegraphics[width=0.95\textwidth]{slides-images/slide-058.jpg}
\caption{3D 生成前沿：从单张图像生成 3D 模型\protect\footnotemark}
\end{figure}
\footnotetext{Slides 第58页。}

\subsection{本章小结}

\begin{itemize}
    \item 早期 3D 生成直接在 3D 数据上训练，受限于数据规模
    \item 可微渲染打通了 3D 生成与 2D 监督之间的桥梁
    \item 利用 2D 扩散模型先验（如 SDS）实现了高质量的文本到 3D 生成
    \item 3D 生成正朝着大规模、多类别、基础模型的方向发展
\end{itemize}

%% ========== 总结与延伸 ==========

